% Root, B and D full literal audits passed; both independent numerical gates passed.
\subsection{Seven source elements and seven states opposite the right-zero two-band}

\begin{lemma}[The exact square-band source partition at order seven]
\label{C:seven-graph-source-algebras}
In the accepted complete list of forty-three order-seven
monoids with positive rational second cohomology, forty-one
have minimal band $2\times2$.  Their contracted algebra
$A=\mathbb QM/\mathbb QK$ and one-dimensional left row
module $V$ have the following exact partition:
\[
\begin{array}{c|r}
V\text{ projective over }A&23\\
A\text{ commutative, hence }V\text{ on a local factor}&12\\
A\cong T_2(\mathbb Q),\quad V\text{ nonprojective}&6.
\end{array}
\]
The six final source indices are $3,10,14,16,19,36$ in
\path{code/outputs/a_seven_cone_targets.json}.
\end{lemma}
\begin{proof}
The complete source construction is the one used in
Lemma~\ref{C:triangular-eleven}: it gives $229$ normalized
tables and $41$ ordinary square-band classes, followed by
the identity-adjoined $2\times3$ and $3\times2$ bands.
The new program
\path{code/a_seven_graph_source_algebras.py}
checks the identity and every associativity triple of all
forty-three literal sources, then reconstructs each
three-dimensional contracted algebra on the outside
labels $0,5,6$ and its actual row character $\chi$.

For each projective case it supplies $z\in A$ with
$az=\chi(a)z$ on every basis element and $\chi(z)=1$.
These equations imply $z^2=z$ and split the row-character
quotient by the projective left ideal $Az$.
The next twelve algebras are checked literally to be
commutative.  A finite-dimensional commutative rational
algebra is a product of Artinian local algebras, and the
one-dimensional character is supported on one local factor
with residue $\mathbb Q$.

For each remaining source the certificate supplies a
basis $e,f,a$ with $e+f=1$, $e^2=e$, $f^2=f$, and
\[
 ef=fe=a^2=fa=ae=0,\qquad ea=af=a,
 \qquad \chi(e)=\chi(a)=0,\quad\chi(f)=1.
\]
Thus these are literal upper-triangular algebras with the
stated nonprojective row character.  All rational entries,
basis ranks and products are checked exactly using
\texttt{Fraction}; no modular-rank inference is used.
The complete certificates are in
\path{code/outputs/a_seven_graph_source_algebras.json}.
The independent checker
\path{code/d_verify_seven_graph_source_algebras.py}
imports no producer or matrix library.  It checks all
$14749$ multiplication triples in the forty-three
sources, derives every row character from its literal
band action, and reproduces all $23+12+6$ cases.
Exact inconsistency checks confirm the nonprojective
cases, and every Peirce product and basis rank is
verified over $\mathbb Q$.  Its complete output is
\path{code/outputs/d_seven_graph_source_algebras_independently_verified.json}.
\end{proof}

\begin{lemma}[The necessary core graph and its three-orbit case]
\label{C:seven-state-two-column-core}
Let $M$ have a $2\times2$ minimal band and $|M|\leq9$.
Suppose that the right-zero-two-band collage on at most
seven states has type $\mathbb CP^2$.  After removing the
already impossible common-fixed-state and common-retraction
cases, the minimal band has exactly two state orbits, each
of size two.  Their core graph has one vertex orbit and
one loop-edge orbit.
\end{lemma}
\begin{proof}
Write $P,Q$ for the two left retractions and $\Gamma$ for
their natural graph.  Let $W=H_1(\Gamma;\mathbb Q)$.
The graph-component exact sequence makes the kernel of
the edge map
\[
 H^2(X;\mathbb Q)\longrightarrow
 \operatorname{Hom}_{\mathbb QM}(W,\mathbb Q)
\]
the image of the ordinary component-act cofiber.
By Lemma~\ref{C:nine-unbounded-act-cups} that image has
zero squares.  A projective-plane generator must therefore
have a nonzero edge functional $\lambda$.
For every band idempotent $e$, invariance gives
$\lambda(w)=\lambda(we)$, so $We\ne0$.

A singleton minimal-band orbit would give a band-fixed
state.  If it is an edge state, either endpoint is a
band-fixed vertex.  A band-fixed vertex is $M$-fixed,
since $vm=(ve)m=v(em)=v$.  This gives zero products.
Thus every orbit in $SK$ has two states, one for each
column.  There are at most three such orbits.

The image graph $\Gamma e$ has one state for each of
these orbits; the left retractions give its vertex and
edge structure.  Its first homology is $We\ne0$.
One orbit is insufficient.  With two orbits the only
possibility is one vertex and one loop edge, as asserted.

Suppose there are three orbits.  A graph with three
states and positive first homology has either one vertex
and two loops, or a loop together with an isolated vertex.
Hence $P=Q$ on all six states of $SK$.
If they agree on every state, the common-retraction
argument of Theorem~\ref{C:right-zero-five-state-reduction}
applies without its state bound: an ordinary-act retract
and a suspension quotient exclude a rank-one nonzero
square.

Otherwise there is exactly one further state $z$, and
$Pz\ne Qz$.  Both are core vertices.  Indeed, the
common-image condition rules out either endpoint being
$z$: membership of $z$ in that image would force
$Pz=Qz=z$.  Commutation with
the right band and equality of $P,Q$ on $SK$ give
$(Pz)k=(Qz)k$ for every $k\in K$.  In a disjoint union
of full two-column orbits this says that the two endpoints
are the two different states of one orbit.

Use the literal quotient complex $D$ from
Theorem~\ref{C:two-column-graph-localization}.
Its only remaining edge is $z$.  The two whole column
core graphs have been killed, so its degree-zero space
is the one-dimensional difference of the two column
vertex classes.  The endpoints of $z$ are in different
columns, giving
\[
 D=[\,\mathbb Q\xrightarrow{\ \pm1\ }\mathbb Q\,].
\]
It is acyclic.  The localization theorem makes
$H_4(X;\mathbb Q)=0$, a contradiction.  This excludes
the three-orbit case.
\end{proof}

\begin{lemma}[The complete two-orbit seven-state control]
\label{C:seven-state-triangular-graph-control}
For the six upper-triangular sources of
Lemma~\ref{C:seven-graph-source-algebras}, no seven-state
right-zero-two-band collage whose minimal band has the
two-orbit graph in Lemma~\ref{C:seven-state-two-column-core}
has type $\mathbb CP^2$.
\end{lemma}
\begin{proof}
Label the two minimal-band orbits by $(0,1)$ and $(2,3)$,
with the first as the vertex orbit and the second as the
loop-edge orbit.  The two left retractions both restrict
to $(0,1,0,1)$ on these four states.
Each of the three remaining states has a row-orbit
profile $\gamma_s\in\{0,1\}^2$.  State relabeling sorts
their multiset.  The actual action of each of the two
outside source elements must send a transient state to
a state with profile $\gamma_s\circ\sigma_m$, and its
action on the two core orbits is the given column map.
Every possible such image is tried.  All four outside
products and then all $49$ source multiplication laws
are checked literally.

For each right action, the program tries every possible
commuting idempotent extending $(0,1,0,1)$, and pairs
exactly those with the same image.  Their transient
values can only be $0,1$ or a transient of profile
$(0,0)$.  This is precisely the necessary and sufficient
condition imposed before the idempotence and full
commutation checks.  Common retractions were already
excluded.  The only identifications used are permutations
of equal-profile transient states and interchanging
$P,Q$.

The complete literal counts are
\[
\begin{array}{c|r|r|r}
\text{source index}&\text{right actions}&
   \text{raw pairs}&\text{retained categories}\\\hline
3&446&2450&1270\\
10&198&2185&1152\\
14&446&2450&1270\\
16&198&1396&792\\
19&1426&3666&1769\\
36&1360&3714&1781.
\end{array}
\]
Thus $4074$ right actions and $15861$ raw pairs give
$8034$ literal seventeen-arrow categories.  Of these,
$4161$ fail the required mod-two second Betti number,
and every one of the other $3873$ has a nonboundary
two-cocycle with an explicitly checked square primitive.
The source and complete certificates are
\path{code/a_seven_graph_T2_acts.py} and
\path{code/outputs/a_seven_graph_T2_acts.json}.

The standalone checker
\path{code/b_verify_seven_graph_T2_acts.py}
uses no producer or matrix-library imports.  It freshly
tries all $7^3$ outside transient assignments, then every
$7^3$ retraction assignment, and expands every retained
orbit under the full stabilizer.  It reproduces all
six counts above and checks $7415382$ typed associativity,
$5495256$ boundary-square and $13756896$ square-primitive
equations.  Its complete source-linked output is
\path{code/outputs/b_seven_graph_T2_acts_independently_verified.json}.
This is a full independent numerical enumeration of the
stated normal form, not merely a sample check.
\end{proof}

\begin{lemma}[The nonsquare seven-element band source]
\label{C:two-by-three-seven-state-graph}
Let $M=1+(I_2\times J_3)$.  No right-zero-two-band
collage on at most seven states has type $\mathbb CP^2$.
\end{lemma}
\begin{proof}
Retain the nonzero invariant edge functional $\lambda$
and the exclusion of band-fixed vertices and common
retractions from the preceding arguments.  Every
minimal-band state orbit is a quotient of its three
columns, so has two or three states.  Its column
partition records exactly when two columns give the
same state.  A commuting map between core orbits can
only coarsen this partition.

For $e=x_{00}$, the graph $\Gamma e$ has one state for
each core orbit and has nonzero first homology.
Thus there are two or three core orbits.
With two, there is one vertex orbit and one loop-edge
orbit; the edge partition refines the vertex partition.

If that edge orbit has three states, use the literal
left trivial resolution
\[
 0\longrightarrow(\mathbb QM e)^2
 \xrightarrow{(b_1,b_2)}
 \mathbb QM(1-e)\xrightarrow{\ {-}\,a\ }
 \mathbb QM e\longrightarrow\mathbb Q\longrightarrow0,
\]
where $a=x_{10}-e$ and $b_j=x_{0j}-e$.
It is exact by the same basis calculation as
Lemma~\ref{C:P5-trivial-two-step}: the four vectors
$x_{ij}-x_{i0}$, $i=0,1$, $j=1,2$, are the kernel
of the middle map.  In particular the trivial module
has projective dimension two.
The space $We$ is spanned by the one core loop $c_0$,
whose two other column images are distinct loop edges
$c_1,c_2$.  The second tensor differential is
\[
 (\alpha c_0,\beta c_0)\longmapsto
       \alpha(c_1-c_0)+\beta(c_2-c_0).
\]
Distinct edge supports make it injective.  Therefore
$\operatorname{Tor}_2^{\mathbb QM}(W,\mathbb Q)=0$.
The component terms in the graph-complex exact sequence
are $\operatorname{Tor}_3$ and $\operatorname{Tor}_4$,
which vanish by this resolution.  Hence $H_4(X;\mathbb Q)=0$.

If both core orbits have two states, they have the same
column partition.  The whole right action then factors
through the column-collapse quotient $q:M\to P_5$,
since every image of a band element lies in these
core orbits.  Choosing a representative of each column
block gives an actual monoid section of $q$.
The corresponding graph cofiber over $P_5$ is an actual
homotopy retract of the cofiber over $M$.
The edge-functional map under this section is the
identity on
$\operatorname{Hom}_{M}(W,\mathbb Q)
=\operatorname{Hom}_{P_5}(W,\mathbb Q)$.
Thus the nonzero $\lambda$ makes the projective-plane
generator restrict nontrivially to that retract.
Its pullback spans $H^2(X;\mathbb Q)$, contradicting
the universal $P_5$ graph cup theorem.

Finally suppose that there are three core orbits.
The core graph has only loops, so $P=Q$ on $SK$.
If any orbit has three states, $|SK|\geq7$ and there
are no remaining states; this is the common-retraction
case.  Otherwise all three orbits have two states and
there is at most one transient.

If the core graph has one vertex orbit and two loop-edge
orbits, all three column partitions coincide.  The same
split $P_5$ quotient applies.  If it has a loop-edge
orbit $B$ at a vertex orbit $A$, and an isolated vertex
orbit $C$, the partitions of $A,B$ agree.  Equal
partitions for $C$ again give the split quotient.
Assume instead that its partition is different.

For $P\ne Q$ there must be one transient $z$.
As before, $Pz,Qz$ are the two different points of
one core vertex orbit.  If that orbit is $A$,
commutation with each row of the band forces
$zK\subseteq A\cup B$.  If it is $C$, commutation
forces $zK\subseteq C$.  In either case the graph's
components supported on $A\cup B$ and those supported
on $C$ remain in two disjoint nonempty invariant
families under $M=1+K$.  The component right act
therefore has at least two orbits, so
\[
 H_1(X;\mathbb Q)
   =\widetilde{\mathbb Q\pi_0\Gamma}
      \otimes_{\mathbb QM}\mathbb Q\ne0.
\]
This excludes the final case.
\end{proof}

\begin{theorem}[Seven source elements and seven states opposite $R_3$]
\label{C:seven-source-seven-state-right-zero}
If $R_3$ is the identity-adjoined two-element right-zero
band, no $(R_3,M)$-collage with $|M|\leq7$ and at most
seven cross states has type $\mathbb CP^2$.
In particular the entire seventeen-arrow profile
$(7,3,7)$ is excluded.
\end{theorem}
\begin{proof}
Sources of order at most six are covered for every
state size by
Corollary~\ref{C:six-source-right-zero-all-states}.
For a seven-element source with zero rational second
cohomology, the two vertex restrictions vanish and
the relative suspension kills the square.
Otherwise use the complete forty-three-source list.

For the forty-one square-band sources, the projective
and local-factor alternatives of
Theorem~\ref{C:two-column-graph-localization} cover the
first $23+12$ cases.  The remaining six reduce by
Lemma~\ref{C:seven-state-two-column-core} to the complete
control in Lemma~\ref{C:seven-state-triangular-graph-control};
fewer than seven states were already excluded by
Theorem~\ref{C:right-zero-six-small-sources}.
The source $1+(I_3\times J_2)$ has contracted algebra
$\mathbb Q$ and projective row module, so is covered
by the same localization theorem.  The other nonsquare
source is Lemma~\ref{C:two-by-three-seven-state-graph}.
These cases exhaust the source list.
\end{proof}
